% Kites Vector Extension (RVV 1.0) - Design Document
% Build: tectonic design_document.tex   (or: make)
\documentclass[11pt]{article}
\usepackage{kitesdoc}
\usepackage[capitalise,noabbrev]{cleveref}
\renewcommand{\kdocshorttitle}{Kites Vector Extension --- Design}
\hypersetup{pdftitle={Kites Vector Extension (RVV 1.0): Design Document}}

\begin{document}

% =====================================================================  TITLE
\begin{titlepage}
\thispagestyle{empty}
\vspace*{1.2cm}
{\sffamily\small\color{kgrey} CS2323 Computer Architecture \quad$\cdot$\quad IIT Hyderabad}\par
\vspace{0.6cm}
{\sffamily\bfseries\fontsize{26}{31}\selectfont\color{kblue} Vector Extension Support\\[2pt] for the Kites RISC-V Simulator\par}
\vspace{0.5cm}
{\sffamily\Large\color{kteal} Design Document\par}
\vspace{0.2cm}
{\large RISC-V ``V'' extension 1.0, the Zve* subsets, and the Zvfh, Zvbb/Zvkb and Zvbc extensions\par}
\vspace{1.4cm}

\begin{tabularx}{\textwidth}{@{}L{3.6cm}Y@{}}
\toprule
\textsf{\textbf{Document}}   & Design document \\
\textsf{\textbf{Version}}    & 1.0 \\
\textsf{\textbf{Author}}     & Raam M \\
\textsf{\textbf{Repository}} & \url{https://github.com/Raam-M/Kites}, branch \texttt{vector-extension} \\
\textsf{\textbf{Companions}} & Scope document (requirements, decisions D1--D14, findings F1--F19); validation report (results) \\
\bottomrule
\end{tabularx}

\vspace{1.0cm}
\begin{keybox}[title=Purpose of this document]
This document describes how the vector extension is built into Kites: the data structures, the
execution model, how each instruction family is implemented, how vector instructions run on the
single-cycle and the pipelined processors, the configuration options, the statistics and timing
model, and the user interfaces. It records the design decisions and the known limitations. The
\emph{why} of the scope is in the scope document; the evidence that the design works is in the
validation report.
\end{keybox}
\vfill
\end{titlepage}

\tableofcontents
\clearpage

% =====================================================================  1
\section{Overview}\label{sec:overview}

Kites now implements the complete RISC-V vector extension~1.0~\cite{spec}: every one of the 627
\textbf{V} instructions (375 base encodings and 252 segment variants), the standard vector
pseudo-instructions, and the seven vector CSRs. On top of that it implements:

\begin{itemize}
\item the embedded subsets Zve64d, Zve64f, Zve64x, Zve32f and Zve32x as selectable profiles;
\item Zvfh (half-precision floating point), Zvbb and its subset Zvkb (bit manipulation) and Zvbc
      (carry-less multiply), 20 further base encodings;
\item VLEN from 128 to 1024 bits and a choice of agnostic-element policy;
\item execution on the single-cycle processor and on all four 5-stage pipeline models;
\item vector statistics and an analytic timing model with lanes and chaining;
\item a vector register viewer, a settings page, vector-unit wires in the circuit view, vector
      commands in the CLI REPL, and vector instructions in custom pseudo-instructions.
\end{itemize}

\Cref{fig:arch} shows the parts and how they connect. The design rests on three ideas:

\begin{enumerate}
\item \textbf{One ISA table.} A table generated from riscv-opcodes~\cite{opcodes} describes every
      vector instruction: name, fixed bits, operand fields and defining extension. The assembler,
      disassembler, executor dispatch, editor support and tests all read it (D10).
\item \textbf{Atomic execution.} A vector instruction reads a snapshot of the state, checks every
      legality rule, computes its results into a staging area and only then commits them in one
      step, recording each change for undo (D5). An illegal instruction therefore changes nothing.
\item \textbf{Processor-independent vector unit.} The unit is a function,
      \texttt{rvv::execute(word, context)}, that sees the register file, the memory controller
      and the step record. The single-cycle processor calls it in its execute step; the pipelines
      call it in write-back (\cref{sec:pipeline}).
\end{enumerate}

\begin{figure}[H]
\centering
\begin{tikzpicture}[
  box/.style={draw=kblue,thick,rounded corners=2pt,fill=klight,align=center,font=\sffamily\small,minimum height=0.95cm,inner sep=4pt},
  new/.style={box,fill=knew},
  arr/.style={-{Stealth[length=2.2mm]},thick,kgrey}]
\node[new] (table) at (0,0) {Vector ISA table\\\scriptsize\texttt{vector\_isa\_table.inc} (riscv-opcodes)};
\node[box] (lexer) at (-4.6,1.9) {Lexer\\\scriptsize vector tokens};
\node[new] (vparse) at (-4.6,0) {Vector parser\\\scriptsize\texttt{v\_formats.cpp}};
\node[box] (prog) at (-4.6,-1.9) {Assembled program};
\node[box] (rvss) at (0,-1.9) {RVSS\\\scriptsize execute step};
\node[box] (rv5s) at (4.6,-1.9) {RV5S pipelines\\\scriptsize write-back};
\node[new] (unit) at (2.3,-3.9) {Vector unit\\\scriptsize\texttt{rvv::execute}};
\node[box] (regs) at (-1.4,-5.9) {Register file\\\scriptsize\texttt{v0}--\texttt{v31}, vector CSRs};
\node[box] (mem) at (2.3,-5.9) {Memory controller\\\scriptsize caches};
\node[new] (prof) at (6.0,-5.9) {Vector profiler\\\scriptsize statistics, timing};
\node[new] (cfg) at (4.6,0) {Vector configuration\\\scriptsize VLEN, profile, extensions};
\node[new] (ui) at (4.6,1.9) {GUI / CLI\\\scriptsize viewer, settings, REPL};
\draw[arr] (lexer) -- (vparse);
\draw[arr] (vparse) -- (prog);
\draw[arr] (table) -- (vparse);
\draw[arr] (prog) -- (rvss);
\draw[arr] (prog.south) |- ++(0,-0.4) -| (rv5s.south west);
\draw[arr] (rvss) -- (unit);
\draw[arr] (rv5s) -- (unit);
\draw[arr] (table) -- (unit);
\draw[arr] (unit) -- (regs);
\draw[arr] (unit) -- (mem);
\draw[arr] (unit) -- (prof);
\draw[arr] (cfg) -- (rv5s.north);
\draw[arr] (ui) -- (cfg);
\draw[arr] (cfg.south east) to[out=-30,in=60] (unit.north east);
\end{tikzpicture}
\caption{Architecture. Shaded yellow: new components; blue: existing components with hooks.}\label{fig:arch}
\end{figure}

% =====================================================================  2
\section{Instruction set table}\label{sec:isa}

\subsection{Generation}
\file{tools/vector/gen_vector_isa.py} reads the riscv-opcodes files \file{rv_v}, \file{rv_zvbb}
and \file{rv_zvbc} and writes \file{src/processor/vector/vector_isa_table.inc}, one line per base
encoding:

\begin{lstlisting}[basicstyle=\ttfamily\scriptsize]
KITES_RVV_INSN("vadd.vv", 0x00000057u, 0xfc00707fu, kVm | kVs2 | kVs1 | kVd, kExtV)
KITES_RVV_INSN("vror.vi", 0x50003057u, 0xf800707fu, kZimm6 | kVm | kVs2 | kVd, kExtZvkb)
\end{lstlisting}

\texttt{match} holds the fixed bits, \texttt{mask} marks them, the fourth argument lists the
variable fields and the last one names the defining extension. The generator checks that the
fixed and variable fields cover all 32 bits exactly and that no instruction word matches two
entries. Zvkb is the subset of Zvbb that riscv-opcodes' \file{rv_zvkn} imports
(\insn{vandn}, \insn{vbrev8}, \insn{vrev8}, \insn{vrol}, \insn{vror}).

\subsection{Run-time table}
\file{vector_isa.cpp} expands the 395 base entries into 647 \texttt{InsnInfo} records: each
unit-stride, strided and indexed load or store with an \texttt{nf} field gets seven segment
variants (\insn{vlseg2e8.v} \ldots{} \insn{vlseg8e8.v}). Each record holds the mnemonic, its base
and suffix, \texttt{match}/\texttt{mask}, the field set, the major opcode and \texttt{funct3}, the
segment count, the extension, and the ordered list of \emph{assembly operands} derived from the
fields (\cref{tab:operands}). Multiply-add instructions list the multiplier before \texttt{vs2},
as the assembly syntax requires.

\begin{table}[H]
\centering\small
\caption{Operand kinds derived from the fields.}\label{tab:operands}
\begin{tabularx}{\textwidth}{@{}L{3.2cm}L{3.6cm}Y@{}}
\toprule
Kind & Fields & Assembly form \\
\midrule
Vector register & \texttt{vd}, \texttt{vs3}, \texttt{vs1}, \texttt{vs2} & \texttt{v0}--\texttt{v31} \\
Integer register & \texttt{rd}, \texttt{rs1}, \texttt{rs2} & \texttt{x0}--\texttt{x31} or ABI names \\
FP register & \texttt{rd}/\texttt{rs1} with OPFVV/OPFVF & \texttt{f0}--\texttt{f31} or ABI names \\
Memory & \texttt{rs1} of a load/store & \texttt{(rs1)} or \texttt{0(rs1)} \\
Immediate & \texttt{simm5}, \texttt{zimm5}, \texttt{zimm6} & $-16..15$, $0..31$, $0..63$ \\
vtype immediate & \texttt{zimm10}, \texttt{zimm11} & \texttt{e32, m2, ta, ma} \\
\texttt{v0} & fixed \texttt{vm}=0 & carry-in or merge selector \\
Optional mask & \texttt{vm} & optional \texttt{, v0.t} \\
\bottomrule
\end{tabularx}
\end{table}

\texttt{zimm6} (\insn{vror.vi}) is split over bit~26 and bits 19..15; \texttt{getField} and
\texttt{setField} handle it as one 6-bit value. \texttt{decode(word)} finds the entry whose
\texttt{mask}/\texttt{match} fits the word. \texttt{isVectorInstruction(word)} recognises the
vector encoding space: all of OP-V, and LOAD-FP/STORE-FP with widths 000, 101, 110 and 111.
Widths 010 and 011 remain \insn{flw}/\insn{fld}/\insn{fsw}/\insn{fsd} (F6).

% =====================================================================  3
\section{Assembler and disassembler}\label{sec:asm}

\paragraph{Lexing.} The lexer produces \texttt{VEC\_REGISTER} for \texttt{v0}--\texttt{v31},
\texttt{VEC\_MASK} for \texttt{v0.t}, and \texttt{VTYPE\_FIELD} for \texttt{e8}\ldots\texttt{e64},
\texttt{mf8}\ldots\texttt{m8}, \texttt{ta}/\texttt{tu}, \texttt{ma}/\texttt{mu} only on lines
whose opcode is \insn{vsetvli} or \insn{vsetivli}. Elsewhere those words remain label references,
so a label called \texttt{ta} still works (AS-04).

\paragraph{Parsing.} \texttt{Parser::parse\_vector} handles every vector mnemonic generically. It
splits the operand tokens into comma-separated groups and matches them against the
instruction's operand list (\file{v_formats.cpp}, \texttt{encodeVector}). Each group must have the
right kind; an immediate is checked against its range; the optional mask may be present only
where \texttt{vm} is a field. The word is built from \texttt{match} with \texttt{setField} and
attached to the instruction unit, so the code generator emits it unchanged. Errors report the
line, the column and the expected syntax from \texttt{expectedSyntax}, for example
``\texttt{vadd.vv vd, vs2, vs1[, v0.t]}''.

\paragraph{Pseudo-instructions.} \texttt{parse\_vector\_pseudo} expands the standard vector
pseudo-instructions (\insn{vneg.v}, \insn{vnot.v}, \insn{vwcvt.x.x.v}, \insn{vncvt.x.x.w},
\insn{vfneg.v}, \insn{vfabs.v}, \insn{vmmv.m}, \insn{vmnot.m}, \insn{vmclr.m}, \insn{vmset.m},
the reversed compares, \insn{vmsge(u).vx} with its masked and temporary-register forms, and
\insn{vl<n>r.v}). \texttt{parse\_csr\_pseudo} adds \insn{csrr}, \insn{csrw}, \insn{csrs},
\insn{csrc} and their immediate forms, with vector CSR names.

\paragraph{Disassembly.} \texttt{rvv::disassemble} prints the mnemonic and the operands in table
order. The disassembly view, the register dump and the error messages use it. Assembling its
output gives the same word (AS-02).

% =====================================================================  4
\section{Architectural state}\label{sec:state}

\texttt{RegisterFile} holds the vector registers as one byte array of $32\times\mathrm{VLENB}$
bytes. Element $i$ of a register occupies the bytes $i\cdot\mathrm{SEW}/8$ upward,
little-endian, so element~0 is in the lowest bytes. A register group of LMUL registers is
contiguous in this array, which makes the group addressing of
\cref{sec:execution} a simple offset calculation. \texttt{ReadVr}/\texttt{WriteVr} copy one
register; \texttt{GetVrBytes} exposes the whole file to the viewer.

\begin{table}[H]
\centering\small
\caption{Vector CSRs.}\label{tab:csrs}
\begin{tabularx}{\textwidth}{@{}L{1.8cm}L{1.6cm}L{1.9cm}Y@{}}
\toprule
CSR & Address & Access & Implementation \\
\midrule
\texttt{vstart} & 0x008 & read/write & reset to 0 after every vector instruction \\
\texttt{vxsat}  & 0x009 & read/write & bit 0 only \\
\texttt{vxrm}   & 0x00A & read/write & bits 1:0 only \\
\texttt{vcsr}   & 0x00F & read/write & composite: reads \texttt{(vxrm << 1) | vxsat}, writes both \\
\texttt{vl}     & 0xC20 & read-only  & written only by \insn{vset*} and fault-only-first loads \\
\texttt{vtype}  & 0xC21 & read-only  & \texttt{vill} in bit 63 \\
\texttt{vlenb}  & 0xC22 & read-only  & VLEN/8, set when the register file is reset \\
\bottomrule
\end{tabularx}
\end{table}

Writing a read-only CSR (address bits 11:10 = 11) is an illegal instruction on every processor.
Reset sets \texttt{vtype.vill}=1 and \texttt{vl}=0 and clears the vector registers (D4). VLEN is
read from the configuration on reset, so a new VLEN applies at the next reset or assemble.

% =====================================================================  5
\section{Execution model}\label{sec:execution}

\subsection{One instruction, one atomic step}

\texttt{rvv::execute(word, context)} runs one instruction. The context holds references to the
register file, the memory controller and the current step record (\texttt{StepDelta}), and
optionally the profiler and the processor cycle. Execution has three phases:

\begin{enumerate}
\item \textbf{Set-up.} An \texttt{ExecState} decodes the fields, reads \texttt{vtype},
      \texttt{vl}, \texttt{vstart} and VLEN, and takes a snapshot of the vector registers
      (\texttt{source}) and a working copy (\texttt{destination}). The instruction's extension
      must be enabled (\cref{sec:config}).
\item \textbf{Checks and computation.} The family handler checks every legality rule first:
      \texttt{vill}, register-group alignment (\texttt{checkGroup}), EMUL within $[1/8, 8]$,
      the overlap rules (\texttt{checkOverlap}, \texttt{forbidOverlap},
      \texttt{checkMaskOverlap}), \texttt{vstart}=0 where required, and the supported SEW. Then
      it loops over the elements. Reads come from the snapshot and writes go to the working copy,
      so an instruction whose source and destination overlap reads the old values.
      Scalar results, CSR updates, \texttt{fflags} and stores are staged.
\item \textbf{Commit.} \texttt{commit} writes the changed vector registers, the staged scalar
      register, the CSRs and the memory, and records every old and new value in the step
      record. It then sets \texttt{vstart} to 0.
\end{enumerate}

A rule violation throws \texttt{IllegalInstruction}, and an access outside memory throws
\texttt{AccessFault}, before phase~3, so nothing has changed (D5). The processor leaves the PC on
the instruction, and Kites shows the message with the source line.

\subsection{Element loop, masking and policies}

For every body element $i$ with $\texttt{vstart} \le i < \texttt{vl}$, the handler asks
\texttt{active(i)}: true when the instruction is unmasked or bit $i$ of \texttt{v0} is set.
Active elements get the result. Inactive elements and tail elements ($i \ge \texttt{vl}$) go
through \texttt{inactive()} and \texttt{tail()}, which apply the policy: undisturbed
leaves the element alone; with the all-ones fill selected (\cref{sec:config}) and
\texttt{vma}/\texttt{vta} set, the element becomes all ones. Mask-producing instructions treat
their tail as agnostic, as the specification requires. Elements below \texttt{vstart} are never
touched.

\subsection{Register groups}

\texttt{group(reg, eew)} returns the register, the element width and
$\mathrm{EMUL} = (\mathrm{EEW}/\mathrm{SEW})\cdot\mathrm{LMUL}$. Element $i$ of a group lives at
byte offset $\texttt{reg}\cdot\mathrm{VLENB} + i\cdot\mathrm{EEW}/8$ of the register array, so
groups need no special indexing. Widening, narrowing, extension, index and mask operands all use
the same accessor with their own EEW.

% =====================================================================  6
\section{Instruction families}\label{sec:families}

\subsection{Configuration}
\insn{vsetvli}, \insn{vsetivli} and \insn{vsetvl} compute
$\mathrm{VLMAX} = \mathrm{LMUL}\cdot\mathrm{VLEN}/\mathrm{SEW}$ and set
$\texttt{vl} = \min(\mathrm{AVL}, \mathrm{VLMAX})$ (D2). \texttt{rs1}=\texttt{x0} with
\texttt{rd}$\neq$\texttt{x0} requests VLMAX; \texttt{rs1}=\texttt{rd}=\texttt{x0} keeps
\texttt{vl} and sets \texttt{vill} if VLMAX would change. A vtype with reserved bits, a reserved
\texttt{vlmul}, SEW above ELEN, or $\mathrm{SEW} > \mathrm{LMUL}\cdot\mathrm{ELEN}$ sets
\texttt{vill}=1 and \texttt{vl}=0 (D4). ELEN depends on the profile, so under Zve32* SEW=64 and
LMUL=1/8 are reserved.

\subsection{Loads and stores}
\file{vector_exec_memory.cpp} handles unit-stride, strided, indexed (ordered and unordered),
segment, fault-only-first, mask and whole-register accesses. Each element is a separate access
through the \texttt{MemoryController} in increasing element (and field) order, so caches see the
traffic and the ordered and unordered indexed forms behave the same (D8). Misaligned element
addresses are allowed. Loads write the working copy; stores are staged and committed with the
rest of the instruction. A fault-only-first load that faults on element $i>0$ trims
\texttt{vl} to $i$; on element 0 it is an access fault. An EEW above ELEN is illegal.

\subsection{Integer and fixed point}
\file{vector_exec_integer.cpp} dispatches on the base mnemonic through tables of element
operations: single-width arithmetic, multiply-add, widening and narrowing, extension, carry and
borrow, compares, merges and moves, reductions, mask logical operations, \insn{vcpop.m},
\insn{vfirst.m}, \insn{vmsbf}/\insn{vmsif}/\insn{vmsof}, \insn{viota}, \insn{vid}, slides,
gathers, compress and whole-register moves. Division by zero and signed overflow give the
RISC-V results. The high halves of 64-bit products use portable 64$\times$64-bit helpers, not
\texttt{\_\_int128} (F9). Fixed-point instructions round with \texttt{vxrm} (rnu, rne, rdn,
rod) and set \texttt{vxsat} on saturation.

\subsection{Floating point}
\file{vector_exec_float.cpp} does not use the scalar FP unit, which deviates from the
specification (F16). For SEW = 32 and 64 it uses host IEEE~754 arithmetic: the rounding mode comes
from \texttt{frm} (\texttt{HostRounding}), the host exception flags are cleared before and read
after each element, and they are accrued into \texttt{fflags} in the RISC-V bit order (D9). On
top of that it applies the RISC-V rules explicitly: canonical NaN results; minimumNumber and
maximumNumber for \insn{vfmin}/\insn{vfmax}; quiet and signalling compares; saturating
float-to-integer conversion; round-to-odd for \insn{vfncvt.rod.f.f.w}; and the 7-bit estimates
\insn{vfrec7}/\insn{vfrsqrt7} from the specification's tables, bit-exact including their special
cases. The code is templated on a format descriptor, so each operation is written once for all
widths.

\subsection{Half precision (Zvfh)}\label{sec:half}
There is no portable half-precision host type, so binary16 is handled in software. A half
operation converts its operands exactly to double, computes the operation in double under the
same rounding mode, and rounds the double result to half with \texttt{roundToHalf}. That
function implements all four host rounding modes and raises inexact, overflow and underflow with
tininess detected \emph{after rounding}, as RISC-V requires.

The double rounding does not change the result. Sums, differences, products and integer
conversions of halves are exact in double. For division and square root, rounding first to a
format with $p' \ge 2p+2$ significand bits and then to $p$ bits gives the correctly rounded
result~\cite{figueroa}; here $p=11$ and $p'=53$. For fused multiply-add, the product of two halves is exact in
double; the sum is rounded once to double and then to half. A difference could only arise if the
two addends were more than 53 bit positions apart and the larger one lay exactly halfway between
two halves. The larger addend is then either a half, which cannot lie halfway, or a product of at
least $2^{29}$, which overflows half anyway. Under the directed modes, both roundings go in the
same direction. A signalling NaN operand raises invalid explicitly, because the software
conversion does not. Widening operations convert halves exactly to single and then use the
single-precision path; narrowing conversions round single to half with \texttt{roundToHalf}.

\subsection{Unordered reductions}
\insn{vfredusum} and \insn{vfwredusum} add in element order, which is one of the orders the
specification allows. When no element is active, the specification also allows adding an
additive identity to \texttt{vs1[0]}. Kites makes the same choice as Spike: a NaN scalar becomes
the canonical NaN, and a signalling NaN raises invalid.

\subsection{Bit manipulation and carry-less multiply}
Zvkb adds \insn{vandn}, \insn{vrol}, \insn{vror} (with the 6-bit \insn{vror.vi} immediate),
\insn{vbrev8} and \insn{vrev8}. Zvbb adds \insn{vbrev}, \insn{vclz}, \insn{vctz}, \insn{vcpop.v}
and the widening shift \insn{vwsll}, which zero-extends \texttt{vs2} and shifts by the low
$\log_2(2\cdot\mathrm{SEW})$ bits of the operand. Zvbc adds \insn{vclmul} and \insn{vclmulh},
defined only for SEW=64. The binary operations reuse the single-width loop; the unary ones have
their own small handler.

% =====================================================================  7
\section{Configuration}\label{sec:config}

\texttt{rvv::VectorConfig} (\file{vector_config.h}) holds the run-time options. The CLI flags,
the GUI settings page and the REPL command \texttt{vconfig} all write the same structure.

\begin{table}[H]
\centering\small
\caption{Configuration options.}\label{tab:config}
\begin{tabularx}{\textwidth}{@{}L{2.6cm}L{3.2cm}Y@{}}
\toprule
Option & Values & Effect \\
\midrule
VLEN & 128, 256, 512, 1024 & register width; applied at the next reset \\
Agnostic fill & undisturbed, ones & value of agnostic tail and inactive elements (D3) \\
Profile & V, Zve64d/f/x, Zve32f/x & ELEN, FP widths, 64-bit high multiplies (\cref{tab:profiles}) \\
Extensions & Zvfh, Zvbb, Zvkb, Zvbc & enable the instructions; Zvbb implies Zvkb \\
Timing lanes & 1, 2, \ldots, 32 & timing model (\cref{sec:timing}) \\
Chaining & on, off & timing model \\
\bottomrule
\end{tabularx}
\end{table}

\begin{table}[H]
\centering\small
\caption{Profiles. A ``--'' means the restriction does not apply because SEW=64 is not available.}\label{tab:profiles}
\begin{tabularx}{\textwidth}{@{}L{2cm}L{1.2cm}L{3.2cm}Y@{}}
\toprule
Profile & ELEN & FP element widths & \insn{vmulh*}, \insn{vsmul} at SEW=64 \\
\midrule
V       & 64 & 32, 64 (+16 with Zvfh) & yes \\
Zve64d  & 64 & 32, 64 (+16 with Zvfh) & no \\
Zve64f  & 64 & 32 (+16 with Zvfh) & no \\
Zve64x  & 64 & none & no \\
Zve32f  & 32 & 32 (+16 with Zvfh) & -- \\
Zve32x  & 32 & none & -- \\
\bottomrule
\end{tabularx}
\end{table}

\texttt{validate} rejects combinations the specification does not allow: Zvfh without an F
profile, and Zvbc with ELEN=32. An instruction of a disabled extension, an FP instruction at an
unsupported width, or an EEW above ELEN is an illegal instruction. The assembler accepts every
instruction regardless of the configuration, as hardware assemblers do with the right flags;
legality is decided when the instruction executes.

% =====================================================================  8
\section{Pipelined execution}\label{sec:pipeline}

\subsection{Execute in write-back}

On the four 5-stage models (no hazard detection or forwarding; hazard detection without
forwarding; forwarding without hazard detection; both), vector instructions and Zicsr CSR
instructions execute atomically in \textbf{WB}. Decode treats them like a bubble with no control
signals; they pass through EX and MEM without effect; \texttt{pipeline\_writeback} calls
\texttt{rvv::execute} (or the CSR routine) for the instruction in MEM/WB before anything else in
the cycle. This keeps the vector unit identical for all processors and gives four properties
without extra logic:

\begin{description}[style=nextline,leftmargin=0.6cm,font=\sffamily\bfseries\color{kblue}]
\item[Scalar operands are current.] Every older instruction has already written back, so the
  vector unit reads its scalar operands from the register file with no forwarding.
\item[Memory order is program order.] A Kites step runs WB before MEM. An older scalar store has
  finished its MEM stage in an earlier cycle. A younger scalar load or store runs its MEM stage in
  the same or a later cycle, after the vector access.
\item[Exceptions are precise.] Write-back runs first in a step. An illegal vector instruction
  stops the step before any younger instruction has written anything. The PC is set to the
  instruction for the error report, and stepping again reports the same error.
\item[Wrong-path instructions never execute.] Branches resolve in MEM and jumps in EX. Both flush
  the younger stages before those instructions reach WB.
\end{description}

\subsection{Hazards}

A vector or CSR instruction never waits for an older instruction. A younger scalar instruction
must wait in two cases, checked in IF/ID after the existing hazard unit
(\texttt{writeback\_dependency\_stall}):

\begin{enumerate}
\item it reads a register that an older vector or CSR instruction still in EX/MEM or MEM/WB
      writes (\texttt{rd} of \insn{vset*}, \insn{vmv.x.s}, \insn{vcpop.m}, \insn{vfirst.m},
      \insn{csrr*}; the FP register of \insn{vfmv.f.s});
\item it is a scalar FP operation, which reads \texttt{frm} and writes \texttt{fflags} in EX,
      while any vector or CSR instruction is in flight.
\end{enumerate}

The stall is re-evaluated every cycle. A value produced in WB cannot be forwarded, so a dependent
instruction waits two cycles on both hazard-detecting models, as for a load without forwarding.
The models without hazard detection do not stall: the program must separate such a producer and
its consumer by two instructions, the same rule as for a load on those models.

\begin{figure}[H]
\centering
\begin{tikzpicture}[x=1.15cm,y=0.62cm,font=\sffamily\scriptsize]
\foreach \c in {1,...,9} \node at (\c,0.9) {\c};
\node[anchor=east] at (0.4,0) {\insn{vmv.x.s a0, v1}};
\node[anchor=east] at (0.4,-1) {\insn{addi a1, a0, 1}};
\foreach \s/\c in {IF/1,ID/2,EX/3,MEM/4} \node[draw=kgrey,fill=klight,minimum width=1.0cm] at (\c,0) {\s};
\node[draw=kblue,fill=knew,minimum width=1.0cm] at (5,0) {WB*};
\foreach \s/\c in {IF/2} \node[draw=kgrey,fill=klight,minimum width=1.0cm] at (\c,-1) {\s};
\foreach \c in {3,4} \node[draw=kmust,fill=white,minimum width=1.0cm] at (\c,-1) {stall};
\foreach \s/\c in {ID/5,EX/6,MEM/7,WB/8} \node[draw=kgrey,fill=klight,minimum width=1.0cm] at (\c,-1) {\s};
\end{tikzpicture}
\caption{A scalar instruction that uses a vector result. WB* is where the vector instruction
executes; WB runs first in a cycle, so ID in cycle 5 reads the new \texttt{a0}.}\label{fig:stall}
\end{figure}

\subsection{CSR instructions}
CSR instructions used to be no-ops on the pipelines. They now execute in WB through the same path
as vector instructions, with the same read-only check as the single-cycle processor, so programs
can set \texttt{vxrm} or read \texttt{vl} and \texttt{vlenb}.

\subsection{Scalar defects fixed along the way}
Running vector programs on the pipelines exposed three existing defects, fixed in separate
commits with regression tests: AUIPC computed a register sum instead of $\mathrm{PC}+\mathrm{imm}$
(so \insn{la} gave wrong addresses on every pipeline) and LUI did not sign-extend; the
forwarding models did not forward the integer source of F/D instructions (load/store base,
\insn{fmv.w.x}, \insn{fcvt.s.w}); and \texttt{Run()} paused at the end-of-program breakpoint
while draining, so \texttt{kites-cli} never finished on a pipeline model.

% =====================================================================  9
\section{Statistics and timing model}\label{sec:timing}

\subsection{Statistics}
After a successful commit the vector unit passes an \texttt{InstructionRecord} to the processor's
\texttt{VectorProfiler}: the instruction, SEW, LMUL, \texttt{vl}, VLMAX, \texttt{vstart}, the
number of body, active and masked-off elements, the bytes loaded and stored, and two 32-bit masks
of the vector registers read and written (collected by the element accessors). The profiler
counts instructions per class (configuration, load, store, integer, fixed point, floating point,
reduction, mask, permutation, bit manipulation) and per mnemonic, the average \texttt{vl} and
$\sum\texttt{vl} / \sum\mathrm{VLMAX}$, active and masked-off elements, bytes moved, and SEW and
LMUL use.

\subsection{Timing model}
The model estimates how long the vector part of a program takes on a decoupled vector unit. The
unit issues vector instructions in program order, at most one per cycle, and not before the cycle
in which the processor reached the instruction. Each unit class (integer ALU, multiply, divide,
FP add, FP multiply, FP divide, load, store, reduction, permute, mask) has one pipelined unit with
$L$ lanes. With start-up latency $\lambda$ for the unit and $n$ body elements, an instruction:

\begin{itemize}
\item occupies its unit for $b = \lceil n/L\rceil$ cycles. Strided and indexed accesses move one
      element per cycle, mask logical operations process $64L$ bits per cycle, and reductions add
      $\log_2 L$ cycles to combine the lanes;
\item starts at $s = \max(\text{issue}, s_\text{prev}+1, \text{unit free}, \text{operands ready})$,
      where an operand is ready at the producer's first result $s_p+\lambda_p$ with chaining, or
      at its last result $s_p+\lambda_p+b_p$ without;
\item finishes at $e = s + \lambda + b$, and its unit is free again at $s + b$.
\end{itemize}

The estimate is $\max e$ minus the issue cycle of the first vector instruction. The default
latencies follow the VMIPS example in Hennessy and Patterson~\cite{hp}: 6 cycles for add, 7 for
multiply, 20 for divide and 12 for load and store. Convoys also follow the textbook: consecutive
instructions form one convoy until one needs a unit already used in the convoy or, without
chaining, reads a register written in it. The number of convoys is the chime count.

\begin{table}[H]
\centering\small
\caption{Worked example, \texttt{vl}=8: \insn{vsetivli} issues at cycle 0, \insn{vle32.v v8} at 1,
\insn{vadd.vv v16, v8, v8} at 2.}\label{tab:timing}
\begin{tabularx}{\textwidth}{@{}L{3.4cm}YYYY@{}}
\toprule
 & 4 lanes, chaining & 4 lanes, no chaining & 1 lane, chaining & 1 lane, no chaining \\
\midrule
load: start, first, last & 1, 13, 15 & 1, 13, 15 & 1, 13, 21 & 1, 13, 21 \\
add: start & 13 & 15 & 13 & 21 \\
estimate (cycles) & 21 & 23 & 27 & 35 \\
convoys & 1 & 2 & 1 & 2 \\
\bottomrule
\end{tabularx}
\end{table}

The model is an estimate for comparing code and configurations, not a cycle-accurate
simulation. It does not model vector-to-scalar register dependencies, memory banks, or the
scalar instructions' own timing.

\subsection{Undo and reset}
Each step record keeps the profiler state before and after a vector instruction. Undo restores
the state from before, redo the state from after, and reset clears it, so the statistics always
match the visible program state.

% =====================================================================  10
\section{User interfaces}\label{sec:ui}

\begin{description}[style=nextline,leftmargin=0.6cm,font=\sffamily\bfseries\color{kblue}]
\item[Register panel] The \emph{Vector} section lists \texttt{v0}--\texttt{v31} in hex; the
  \emph{Vector CSRs} section shows \texttt{vl}, the decoded \texttt{vtype}, \texttt{vstart},
  \texttt{vxrm}, \texttt{vxsat}, \texttt{vcsr} and \texttt{vlenb}.
\item[Vector tab] A table of the 32 registers split into elements at the current SEW or a chosen
  one, shown as hex, unsigned, signed or floating point (including half). Tail elements are
  dimmed, masked-off elements can be shaded, LMUL groups are banded, and elements changed since
  the last update are highlighted. The vector CSRs are shown above the table and the statistics
  and timing report below it. Updates are coalesced so a fast run does not flood the GUI thread.
\item[Settings] \emph{Settings $\rightarrow$ Vector Extension} edits the options of
  \cref{tab:config}. They are saved with \texttt{QSettings} and applied at start-up before the
  processors are created.
\item[Circuit view] Every processor design has a \emph{Vector Unit} and a \emph{Vector Regs}
  block with five wires: scalar operands and results to and from the register file, vector
  register reads and writes, and memory. The pipelines have a sixth wire from MEM/WB. The wires
  light up while a vector instruction executes (in WB on the pipelines).
  \file{tools/vector/add_vector_unit_to_circuits.py} adds the blocks without reformatting the
  design files.
\item[CLI] The \texttt{kites-cli} options \texttt{--vlen}, \texttt{--vector-agnostic} and
  \texttt{--vector-ext} set the configuration, and a batch run also writes
  \file{vector_stats.json}. The REPL adds
  \texttt{vregs}/\texttt{vreg} (elements in any SEW and format), \texttt{vcsr},
  \texttt{vconfig} and \texttt{vstats}.
\item[Custom pseudo-instructions] An expansion line may use vector instructions. Each operand is
  checked against the instruction's operand list: vector, integer and FP registers, a
  \texttt{(base)} memory operand, immediate ranges, vtype fields and \texttt{v0.t}. Arguments
  are substituted inside \texttt{(\ldots)} too.
\item[Editor] The editor highlights vector mnemonics and registers and offers them in
  autocompletion; both read \file{resources/instructions/instructions.json}. The profiler counts
  vector instructions as \texttt{V-Type}.
\end{description}

% =====================================================================  11
\section{Design decisions}\label{sec:decisions}

The decisions D1--D14 of the scope document hold, with these refinements:

\small
\begin{longtable}{@{}L{0.9cm}L{6.1cm}L{8.1cm}@{}}
\caption{Decisions added or revised during the implementation.}\label{tab:decisions}\\
\toprule
ID & Decision & Rationale \\
\midrule
\endfirsthead
\caption[]{Decisions (continued).}\\
\toprule
ID & Decision & Rationale \\
\midrule
\endhead
D12$'$ & Supersedes D12: the pipelines execute vector and CSR instructions in WB, with a stall for dependent scalar instructions (\cref{sec:pipeline}). & One vector unit for every processor; precise exceptions and program-order memory access follow from the existing step order. \\
D14$'$ & Supersedes D14: every vector mnemonic assembles; legality depends on the configuration at execution time. & Configuration (profile, extensions) can change without re-assembling; matches how hardware treats unsupported encodings. \\
D15 & Zve* profiles, Zvfh, Zvbb/Zvkb and Zvbc are options; the default configuration is plain V with VLEN=128. & The defaults keep the behaviour of the scope; each option is validated against the specification's combination rules. \\
D16 & Half precision is computed through double with a software rounding step. & Portable (no compiler half type needed), correctly rounded (\cref{sec:half}), and RISC-V tininess detection. \\
D17 & \insn{vfredusum}/\insn{vfwredusum} with no active element canonicalise a NaN scalar. & An allowed choice that matches Spike, the reference model. \\
D18 & The timing model is analytic and separate from the processors' cycle counts. & Shows the effect of lanes and chaining without changing the existing cycle semantics of the single-cycle and pipeline models. \\
D19 & Statistics are part of the undoable state. & Undo, redo and reset keep statistics consistent with the program state. \\
\bottomrule
\end{longtable}
\normalsize

% =====================================================================  12
\section{Known limitations}\label{sec:limits}

\begin{itemize}
\item \texttt{frm}=RMM rounds ties to even in floating-point arithmetic, because the host has no
      such mode. Conversions to integer honour RMM.
\item For SEW 32 and 64, underflow uses the host's tininess rule. Half precision detects
      tininess after rounding.
\item Scalar single and half values are kept zero-extended in the \texttt{f} registers, as in
      Kites' F extension (F18), not NaN-boxed. \texttt{.vf} operands use the low SEW bits.
\item Kites has no traps: an illegal instruction or access fault stops the program, and
      \texttt{vstart} is non-zero only after an explicit CSR write.
\item On the pipeline models without hazard detection, a scalar instruction that reads a vector
      or CSR result must be two instructions after it.
\item The timing model is not cycle-accurate (\cref{sec:timing}).
\end{itemize}

% =====================================================================  13
\section{Source files}\label{sec:files}

\small
\begin{longtable}{@{}L{6.4cm}L{8.7cm}@{}}
\caption{New and changed files.}\label{tab:files}\\
\toprule
File & Role \\
\midrule
\endfirsthead
\caption[]{New and changed files (continued).}\\
\toprule
File & Role \\
\midrule
\endhead
\file{src/processor/vector/vector_isa_table.inc} & generated encodings \\
\file{src/processor/vector/vector_isa.*} & ISA table, decode, operands, disassembly \\
\file{src/processor/vector/vtype.*} & vtype decode/encode, VLMAX, legality \\
\file{src/processor/vector/vector_config.*} & VLEN, policy, profile, extensions \\
\file{src/processor/vector/vector_unit.*} & \texttt{execute}, \texttt{ExecState}, legality helpers, commit \\
\file{src/processor/vector/vector_exec.h} & shared declarations of the families \\
\file{src/processor/vector/vector_exec_integer.cpp} & integer, fixed point, mask, reduction, permutation, Zvbb/Zvbc \\
\file{src/processor/vector/vector_exec_memory.cpp} & loads and stores \\
\file{src/processor/vector/vector_exec_float.cpp} & floating point, half precision \\
\file{src/processor/vector/vector_profiler.*} & statistics and timing model \\
\file{src/processor/vector/vector_format.*} & element formatting for the CLI and GUI \\
\file{src/assembler/parse_formats/v_formats.cpp} & vector parser, pseudo-instructions, CSR pseudos \\
\file{src/assembler/lexer.cpp}, \file{parser.*}, \file{code_generator.*} & vector tokens, dispatch, pre-encoded words \\
\file{src/processor/registers.*} & vector registers and CSRs \\
\file{src/processor/rvss/rvss_processor.*} & vector execution, undo/redo, CSR read-only check \\
\file{src/processor/rv5s/*} & WB execution, stall, CSRs, scalar fixes \\
\file{src/processor/processor_base.*} & step record, profiler, circuit wires \\
\file{src/ui/vector_tab/*}, \file{src/ui/dialogs/settings_dialog.*} & vector viewer, settings page \\
\file{src/ui/register_table/*} & Vector and Vector CSRs sections \\
\file{src/custom_pseudo_manager/*} & vector expansions \\
\file{cli/main.cpp}, \file{cli/repl.cpp} & flags and REPL commands \\
\file{resources/circuit_designs/*.json} & vector unit blocks and wires \\
\file{tools/vector/*} & table and golden generators, Spike harness, circuit script \\
\file{test/vector/*}, \file{test/rv5s_regression_test.cpp} & tests \\
\bottomrule
\end{longtable}
\normalsize

% =====================================================================  REFERENCES
\begin{thebibliography}{99}\small\markboth{References}{References}
\bibitem{spec} RISC-V International. \emph{The RISC-V Instruction Set Manual, Volume I: Unprivileged
  Architecture}, ``V'' Standard Extension for Vector Operations, version 1.0, and the Zvfh, Zvbb,
  Zvkb and Zvbc extensions. \url{https://docs.riscv.org/reference/isa/v20260120/unpriv/v-st-ext.html}
\bibitem{opcodes} RISC-V International. \emph{riscv-opcodes}. \url{https://github.com/riscv/riscv-opcodes}
\bibitem{spike} RISC-V Software. \emph{Spike RISC-V ISA Simulator}.
  \url{https://github.com/riscv-software-src/riscv-isa-sim}
\bibitem{hp} J.~L.~Hennessy and D.~A.~Patterson. \emph{Computer Architecture: A Quantitative
  Approach}, 6th ed., Chapter 4. Morgan Kaufmann, 2017.
\bibitem{figueroa} S.~A.~Figueroa. ``When is double rounding innocuous?'' \emph{ACM SIGNUM
  Newsletter} 30(3), 1995.
\end{thebibliography}

\end{document}
